\section{Methods}
\label{sec:methods}

Every simulated case is one realization of a front-end: its components drawn within
manufacturing tolerances, at most one injected fault, and three electrodes drawn from a
population. Two simulations are run per case. The first takes the measurements that the
instrument could take on itself, with the patient's electrodes connected, and gives the
features. The second measures the specifications with the test networks of the standard and
gives the labels.

\subsection{Circuits}
\label{sec:circuits}

\begin{figure*}[!t]
\centering
\includegraphics[width=\textwidth]{schematic_integrated}
\caption{\cA: integrated instrumentation amplifier with its discrete network, on a single
\SI{3.3}{\volt} supply. $V_\mathrm{cal}$, $V_\mathrm{cmt}$ and $I_\mathrm{lo}$ are the internal
self-test sources; nets with the same name are connected. The electrode model at the lower right
is placed between the body and each of the three leads.}
\label{fig:schematic}
\end{figure*}

Both circuits are single-lead front-ends with the same chain: input protection and
radio-frequency filtering, bias resistors, an instrumentation amplifier, a first-order high-pass
filter, a gain stage, a Sallen-Key low-pass filter and a driven-right-leg (DRL) integrator.

\cA\ (Fig.~\ref{fig:schematic}) is built around a Texas Instruments INA333~\cite{ina333} on a
single \SI{3.3}{\volt} supply; it has 16 resistors, 8 capacitors and 5 discrete op-amps, and a
gain of 278. Its values follow from the standard: a first-stage gain of 4.03 keeps a
$\pm$\SI{300}{\milli\volt} electrode offset within the output swing; the high-pass corner is
\SI{0.041}{\hertz}, as the impulse test requires; the low-pass filter is a Butterworth response
at \SI{185}{\hertz}. \cB\ replaces the integrated amplifier by three op-amps and seven resistors
on $\pm$\SI{5}{\volt} supplies (20 resistors, 6 capacitors, 6 op-amps, gain 414). It has the
resistor symmetries identified in benchmark circuits~\cite{chen2025}; its schematic is in the
supplementary material.

The circuits are simulated with ngspice~\cite{ngspice}. The INA333 is a behavioral model built
from its data sheet (three-amplifier structure, single pole, noise, output clamping, offset and
finite common-mode rejection ratio, CMRR), and the op-amps are single-pole models with offset,
noise and clamping. Against the manufacturer's macromodel of the amplifier alone, gain and CMRR
at \SI{50}{\hertz} agree (\SI{102}{\decibel} against \SI{100}{\decibel}); the macromodel was not
used because its fixed parameters allow neither Monte Carlo variation nor faults in the block.
Healthy variation draws resistors uniformly within $\pm$\SI{1}{\percent} and capacitors within
$\pm$\SI{5}{\percent}, and amplifier parameters within their data-sheet ranges.

\subsection{Electrodes}
\label{sec:electrodes}

Each electrode is a half-cell potential in series with $R_s$ and with $R_d \parallel C_d$. Half
of the cases use gel electrodes, with the network of the standard (\SI{51}{\kilo\ohm} in parallel
with \SI{47}{\nano\farad}). The other half use one of six dry materials (conductive polymer,
conductive fabric with and without a membrane, stainless steel, silver, platinum) with the
parameters fitted in~\cite{joutsen2024}; for each case one of the six subjects of the open data
of that work is drawn and that subject's contact resistance is used, which spans
\SI{0.26}{\mega\ohm} to \SI{68}{\mega\ohm} for the polymer and the fabrics (referred to below as
porous electrodes) and \SI{0.6}{\kilo\ohm} to \SI{0.58}{\mega\ohm} for the solid metals. Every
parameter of each of the three electrodes is then drawn log-uniformly between half and twice its
value, so that imbalance is part of the normal variation. The patient is coupled to mains and
earth as in~\cite{winter1983}. The absence of a subject spread for gel and the factor of two
between electrodes are assumptions.

\subsection{Specifications and Labels}
\label{sec:specs}

Eleven quantities are evaluated for every case following clause 201.12.4 of
IEC~60601-2-25~\cite{iec60601225} (Table~\ref{tab:specs}), with the test networks of the standard
at the inputs and not the patient's electrodes. A case is \emph{compliant} when all are within
their limits. Both nominal circuits and 500 healthy realizations of each are compliant.

\begin{table}[!t]
\caption{Specifications Evaluated for Every Case (IEC~60601-2-25, Clause 201.12.4) and Values of
the Nominal Circuits}
\label{tab:specs}
\centering
\setlength{\tabcolsep}{4pt}
\begin{tabular}{@{}lccc@{}}
\toprule
Specification & Limit & \cA & \cB \\
\midrule
Gain error at \SI{10}{\hertz} & $\le$ \SI{5}{\percent} & 0 & 0 \\
Response, \SIrange{0.67}{40}{\hertz}, deviation & $\le$ \SI{10}{\percent} & \SI{0.3}{\percent} & \SI{0.3}{\percent} \\
Response, \SIrange{40}{150}{\hertz}, minimum & $\ge$ 0.70 & 0.83 & 0.83 \\
Response, \SIrange{40}{500}{\hertz}, maximum & $\le$ 1.10 & 1.00 & 1.00 \\
Baseline shift after impulse$^{\mathrm{a}}$ & $\le$ \SI{100}{\micro\volt} & \SI{75}{\micro\volt} & \SI{75}{\micro\volt} \\
Baseline slope after impulse$^{\mathrm{a}}$ & $\le$ \SI{300}{\micro\volt\per\second} & \SI{18}{\micro\volt\per\second} & \SI{18}{\micro\volt\per\second} \\
Common-mode rejection & $\ge$ \SI{89}{\decibel} & \SI{122}{\decibel} & \SI{111}{\decibel} \\
Noise, peak to valley$^{\mathrm{b}}$ & $\le$ \SI{30}{\micro\volt} & \SI{5.4}{\micro\volt} & \SI{3.4}{\micro\volt} \\
Loss with \SI{620}{\kilo\ohm} $\parallel$ \SI{4.7}{\nano\farad} & $\le$ \SI{20}{\percent} & \SI{3.0}{\percent} & \SI{3.0}{\percent} \\
Gain change with $\pm$\SI{300}{\milli\volt} offset & $\le$ \SI{5}{\percent} & 0 & 0 \\
Input range & $\ge$ \SI{5}{\milli\volt} & \SI{5.8}{\milli\volt} & \SI{6.0}{\milli\volt} \\
\bottomrule
\multicolumn{4}{@{}l}{\footnotesize $^{\mathrm{a}}$\SI{3}{\milli\volt}, \SI{100}{\milli\second}
impulse; referred to the input.}\\
\multicolumn{4}{@{}l}{\footnotesize $^{\mathrm{b}}$Referred to the input, \SIrange{0.05}{150}{\hertz};
taken as 6.6 times the rms value.}
\end{tabular}
\end{table}


Each case carries three labels: the \emph{functional} label (compliance, and the specification
values), the \emph{location} label (the component altered) and the \emph{origin} label (none,
circuit or electrode). Since the specifications belong to the circuit, an electrode fault leaves
the case compliant and is recorded by its origin. The \emph{percentage} label of the diagnosis
literature is kept for comparison: a case is faulty when a component was altered.

\subsection{Fault Model}

One fault at most is injected per case (Table~\ref{tab:faults}). Small parametric deviations are
included on purpose, since that is where functional and percentage severity disagree most. Each
fault condition is simulated 200 times and the healthy circuit \num{5000} times, which gives
\num{63600} cases of \cA\ and \num{66400} of \cB.

\begin{table}[!t]
\caption{Fault Model}
\label{tab:faults}
\centering
\setlength{\tabcolsep}{3pt}
\begin{tabular}{@{}p{0.27\columnwidth}p{0.53\columnwidth}cc@{}}
\toprule
Fault & Model & \multicolumn{2}{c}{Conditions} \\
 & & A & B \\
\midrule
Open, short & \SI{1}{\giga\ohm} in series; \SI{1}{\ohm} in parallel & 48 & 52 \\
Parametric & $\pm$5, $\pm$10, $\pm$20, $\pm$\SI{50}{\percent} of nominal & 192 & 208 \\
Capacitor degradation & $-$\SI{20}{\percent} with \SI{10}{\ohm} in series; $-$\SI{50}{\percent} with \SI{100}{\ohm} & 16 & 12 \\
Op-amp & offset of 20 or \SI{50}{\milli\volt}; open-loop gain divided by 100 or 1000 & 20 & 24 \\
Integrated amplifier & offset of 5 or \SI{20}{\milli\volt}; CMRR of 70 or \SI{50}{\decibel}; gain error of $-$5 or $-$\SI{20}{\percent} & 6 & -- \\
Electrode detached & \SI{1}{\giga\ohm} in series (each of three electrodes) & 3 & 3 \\
Electrode degraded & $R_d$ multiplied and $C_d$ divided by 5 or 20 (each electrode, and both input electrodes) & 8 & 8 \\
\bottomrule
\end{tabular}
\end{table}


\subsection{Self-Test Measurements}
\label{sec:measurements}

Three internal test sources are assumed (Fig.~\ref{fig:schematic}): $V_\mathrm{cal}$ in series
with one lead, for a calibration pulse and a differential tone; $V_\mathrm{cmt}$ at the reference
of the DRL amplifier, which the body follows, for a common-mode tone; and a differential current
$I_\mathrm{lo}$, as in ac lead-off detection. They are ideal sources. The responses are read at
the output by the converter of the instrument, with the patient connected and no signal present.
Four sets of measurements give 26 features:
\begin{itemize}
\item C1, dc levels at the output and at the instrumentation-amplifier output (2 features,
\SI{0.13}{\second});
\item C2, gain, phase and CMRR estimate from differential and common-mode tones at 0.05, 0.5, 10,
50 and \SI{150}{\hertz} (15 features, \SI{94}{\second});
\item C3, seven descriptors of the response to a \SI{1}{\milli\volt}, \SI{200}{\milli\second}
calibration pulse (\SI{1}{\second});
\item C4, response to a \SI{1}{\nano\ampere} lead-off current at 10 and \SI{100}{\hertz}
(2 features, \SI{2}{\second}).
\end{itemize}
A measurement model is applied to the noise-free simulation results: \SI{1}{\milli\volt} rms of
noise at the converter input, 12 bits at \SI{1}{\kilo\hertz}, clipping, dc readings averaged over
64 samples and tones estimated from \num{1000} samples or two periods. Referred to the input of
\cA, one code is \SI{2.9}{\micro\volt}.

\subsection{Testability Analysis}
\label{sec:testability}

What the measurements can distinguish is analyzed before any model is trained. The sensitivity of
feature $f_i$ to component $p_j$ of the nominal circuit is expressed in units of the spread
$\sigma_i$ of that feature over healthy circuits,
\begin{equation}
z_{ij} = \frac{0.01\,p_j}{\sigma_i}\,\frac{\partial f_i}{\partial p_j}.
\label{eq:sensitivity}
\end{equation}
Two components whose columns of $Z=[z_{ij}]$ are parallel act on the measurements in the same
way. Components whose columns have an absolute cosine of at least 0.99 are joined in an \emph{a
priori ambiguity group}; the other components and the amplifiers are groups of one. The number
of singular values $s$ of $Z$ with $10\,s\ge 3$ is reported as the testability rank. In addition,
on the simulated faults, two fault conditions are called indistinguishable when no single feature
separates their medians by three pooled spreads (interquartile range over 1.349); this univariate
criterion is conservative and gives a map against which the models are compared.

\subsection{Models and Evaluation}
\label{sec:models}

\emph{Compliance decision.} Four methods use the same features. The \emph{limit test} fails a
case if any feature leaves the interval of the healthy circuits (\SI{99}{\percent} joint
coverage). The \emph{specification regression} follows alternate test: one gradient-boosting
regressor per specification, and a case fails if any prediction is beyond its limit. The
\emph{classifier} is a gradient-boosting classifier of compliance; with an \emph{escape target}
its threshold is lowered to the value that, on validation data, lets through \SI{1}{\percent} of
the non-compliant cases.

\emph{Severity.} The classifier, with class weights, is trained once on the functional and once
on the percentage label, and both are scored against compliance.

\emph{Location and origin.} On non-compliant cases, five classifiers ($k$-nearest neighbors,
random forest, support vector machine, gradient boosting, multilayer perceptron) name the a
priori group and, separately, the component. Origin is a three-class problem defined by what has
to be done: \emph{none} (healthy, or a harmless deviation), \emph{circuit} (out of specification)
and \emph{electrode} (detached or degraded, or an open protection resistor in series with a lead,
which is electrically the same as a detached electrode). Models are those of
scikit-learn~\cite{pedregosa2011} with small hyperparameter grids.

\emph{Terminology.} In the vocabulary of metrology~\cite{vim,gum}, the specifications are the
measurands and the self-test estimates indirectly whether they are within their limits; the
decision is a conformity assessment in the sense of JCGM~106~\cite{jcgm106}. The \emph{escape
rate} is the fraction of non-conforming circuits accepted and the \emph{false-reject rate} the
fraction of conforming ones rejected. They correspond to the consumer's and the producer's risk
conditioned on the true state of the circuit, over a simulated population whose proportion of
non-conforming cases is set by the fault model and is not that of instruments in service; the
escape target acts as a guard band. Classification accuracy, balanced accuracy, recall and F1 are
figures of merit of classifiers, not statements about measurement accuracy.

\emph{Protocol.} Every experiment is repeated on three random splits (52.5, 17.5 and
\SI{30}{\percent} for training, validation and test), stratified by fault condition, with
thresholds and hyperparameters chosen on validation only. Results are means over repetitions with
\SI{95}{\percent} confidence intervals (Student's $t$), on all electrode types and the 26
features unless stated. Robustness is examined by holding parametric magnitudes out of training,
sweeping converter noise (0.1 to \SI{10}{\milli\volt}) and resolution (8 to 16 bits), scoring on
datasets generated with other tolerances, and changing the common-mode tone. A minimal set of
measurements is found by greedy forward selection over the self-test actions.
